\documentclass[11pt,a4paper]{article}

% === 中文支持 ===
\usepackage{ctex}
\usepackage[T1]{fontenc}

% === 页面布局 ===
\usepackage[top=2.5cm, bottom=2.5cm, left=2.5cm, right=2.5cm]{geometry}

% === 表格 ===
\usepackage{booktabs}
\usepackage{longtable}
\usepackage{array}
\usepackage{enumitem}
\usepackage{tabularx}

% === 数学符号 ===
\usepackage{amssymb}

% === 颜色 ===
\usepackage{xcolor}
\definecolor{primary}{HTML}{1a1a2e}
\definecolor{accent}{HTML}{3b82f6}
\definecolor{success}{HTML}{22c55e}
\definecolor{danger}{HTML}{ef4444}
\definecolor{warning}{HTML}{f59e0b}
\definecolor{graytext}{HTML}{6b7280}

% === 超链接 ===
\usepackage[hidelinks]{hyperref}
\hypersetup{colorlinks=true, linkcolor=primary, urlcolor=accent}

% === 标题样式 ===
\usepackage{titlesec}
\titleformat{\section}{\Large\bfseries\color{primary}}{}{0em}{}[\titlerule]
\titleformat{\subsection}{\large\bfseries\color{primary}}{}{0em}{}
\titleformat{\subsubsection}{\normalsize\bfseries\color{primary}}{}{0em}{}

% === 页眉页脚 ===
\usepackage{fancyhdr}
\pagestyle{fancy}
\fancyhf{}
\fancyhead[L]{\small\color{graytext} 企业级 AI Agent 应用商店 -- V0 需求文档}
\fancyhead[R]{\small\color{graytext} v0.1}
\fancyfoot[C]{\small\color{graytext} \thepage}
\renewcommand{\headrulewidth}{0.4pt}

% === 文档信息 ===
\title{
\vspace{-1em}
{\Huge\bfseries\color{primary} 企业级 AI Agent 应用商店}\\[0.8em]
{\Large\bfseries\color{accent} V0 先备阶段需求文档}\\[0.5em]
{\large 智能体分工、技术选型与文档体系搭建}
}
\author{
\textbf{PM AI}：万能产品小助手（OpenClaw）\\[0.3em]
\textbf{审阅}：魏子政
}
\date{2026-05-06 \quad 版本 v0.1}

\begin{document}

\maketitle
\thispagestyle{empty}

\begin{center}
\begin{tabular}{>{\bfseries}p{3cm} p{10cm}}
\toprule
文档阶段 & V0（先备阶段） \\
文档目标 & 搭建智能体分工体系、锁定技术选型、建立文档体系，为 V1 正式开发做好准备 \\
当前状态 & \textcolor{success}{V0 全部需求已完成} \\
PM AI & 万能产品小助手（OpenClaw） \\
开发者 & Cursor Agent \\
技术栈 & Vue 3 + Element Plus + FastAPI + PostgreSQL + MinIO + OpenSearch + JWT + Docker Compose \\
\bottomrule
\end{tabular}
\end{center}

\vspace{1.5em}

\begin{abstract}
本文档是 V0 先备阶段的完整需求文档，记录了从项目启动到正式开发前的所有需求、决策和交付物。

V0 阶段的核心目标不是写业务代码，而是\textbf{搭好框架}：确定两种智能体怎么配合、技术栈选什么、文档怎么管、需求怎么流转。这些决定会影响后续所有阶段的开发效率和质量。

V0 共包含 7 项需求，全部已完成。V1 阶段将在此基础上接入第一批正式业务需求。
\end{abstract}

\newpage
\tableofcontents
\newpage

% ============================================================
\section{V0 阶段总览}
% ============================================================

\subsection{阶段定义}

V0 是先备阶段，目标是搭建项目基础设施和协作框架，为 V1 正式业务需求开发做好准备。

\begin{table}[h]
\centering
\begin{tabular}{clp{6.5cm}}
\toprule
\textbf{阶段} & \textbf{名称} & \textbf{核心目标} \\
\midrule
V0 & 先备阶段 & 智能体分工 + 技术选型 + 文档体系 + 自闭环机制 \\
V1 & 正式需求 & 第一批业务需求接入（CRUD、搜索、审批、安全扫描） \\
V1.x & 功能迭代 & 功能增强、Bug 修复、性能优化 \\
V2+ & 演进阶段 & 架构升级、生态扩展 \\
\bottomrule
\end{tabular}
\end{table}

\subsection{V0 需求清单}

\begin{table}[h]
\centering
\begin{tabular}{clp{4.5cm}c}
\toprule
\textbf{编号} & \textbf{需求名称} & \textbf{核心内容} & \textbf{状态} \\
\midrule
D-001 & 技术选型锁定 & 确定前后端、数据库、搜索、存储、认证、容器化方案 & \textcolor{success}{已完成} \\
D-002 & 智能体分工定义 & OpenClaw 与 Cursor 的职责划分与协作方式 & \textcolor{success}{已完成} \\
D-003 & Skill/规则体系搭建 & 两套智能体的上下文加载机制与文件清单 & \textcolor{success}{已完成} \\
D-004 & 文档体系与规范 & tex-pdf 格式、命名规范、版本管理、目录结构 & \textcolor{success}{已完成} \\
D-005 & 需求变更流程 & 需求新增、评估、拆解、分配、验收的完整流程 & \textcolor{success}{已完成} \\
D-006 & 自闭环机制 & 5 个闭环定义与监控标准 & \textcolor{success}{已完成} \\
D-007 & 开源方案调研 & 参考外部应用商店架构，蒸馏可借鉴的设计思路 & \textcolor{success}{已完成} \\
\bottomrule
\end{tabular}
\end{table}

% ============================================================
\section{D-001：技术选型锁定}
% ============================================================

\subsection{需求描述}

确定项目全部技术栈，锁定后不得自行变更。技术选型需覆盖前端、后端、数据库、对象存储、搜索引擎、认证授权、容器化、数据库迁移 8 个层面。

\subsection{决策结果}

\begin{table}[h]
\centering
\begin{tabular}{llp{5.5cm}}
\toprule
\textbf{层} & \textbf{技术} & \textbf{选型理由（摘要）} \\
\midrule
前端 & Vue 3 + Element Plus & 学习曲线低、组件库成熟、国内主流 \\
后端 & FastAPI & 异步原生、自动文档、Python 生态 \\
数据库 & PostgreSQL 16 & ACID、JSONB、扩展性强 \\
对象存储 & MinIO & S3 兼容、私有化部署 \\
搜索 & OpenSearch 2.x & 全文+向量一体化、Apache 2.0 \\
认证 & JWT + RBAC & 轻量够用、无状态 \\
容器化 & Docker Compose & MVP 够用、开发体验好 \\
迁移 & Alembic & SQLAlchemy 官方工具 \\
\bottomrule
\end{tabular}
\end{table}

\textbf{硬约束：前端不用 Next.js/React，后端不用 Spring/Java。}

\subsection{精简记录（v1.0 $\to$ v2.0）}

\begin{table}[h]
\centering
\begin{tabular}{llp{4.5cm}}
\toprule
\textbf{原组件} & \textbf{替代} & \textbf{原因} \\
\midrule
Next.js & Vue 3 & 不需要 SSR \\
MongoDB & PG JSONB & 减少组件 \\
Keycloak & JWT + RBAC & 初期不需要完整 IdP \\
Temporal & 数据库状态机 & 审批流程简单 \\
Kafka & asyncio & 异步量不大 \\
K8s & Docker Compose & MVP 单机够用 \\
\bottomrule
\end{tabular}
\end{table}

\subsection{交付物}

\begin{itemize}[nosep]
  \item \texttt{tech-selection.md} -- 技术选型说明书 v2.0
  \item \texttt{.cursor/rules/project-overview.mdc} -- Cursor 技术栈锁定规则
  \item \texttt{CLAUDE.md} -- Cursor 总入口（含技术栈约束）
\end{itemize}

\subsection{验收结论}

\textcolor{success}{通过。} tech-selection.md / CLAUDE.md / project-overview.mdc 三者技术栈一致，无矛盾。

% ============================================================
\section{D-002：智能体分工定义}
% ============================================================

\subsection{需求描述}

明确 OpenClaw（PM AI）和 Cursor Agent 各自的职责范围，消除职责交叉和模糊地带。魏子政明确提出：

\begin{quote}
``你主要负责需求生成拆解和验收。Cursor 负责前后端（包含 UI 设计）、环境、运维、测试（和验收不一样！）。''
\end{quote}

\subsection{决策结果}

\begin{table}[h]
\centering
\begin{tabular}{lcc}
\toprule
\textbf{职责} & \textbf{OpenClaw (PM)} & \textbf{Cursor Agent} \\
\midrule
需求生成与拆解 & \textcolor{success}{\checkmark} & \\
独立验收 & \textcolor{success}{\checkmark} & \\
文档维护与汇报 & \textcolor{success}{\checkmark} & \\
进度管理 & \textcolor{success}{\checkmark} & \\
UI 设计 & & \textcolor{success}{\checkmark} \\
前端开发 & & \textcolor{success}{\checkmark} \\
后端开发 & & \textcolor{success}{\checkmark} \\
环境搭建与运维 & & \textcolor{success}{\checkmark} \\
自验收 & & \textcolor{success}{\checkmark} \\
响应 PM 验收 & & \textcolor{success}{\checkmark} \\
\bottomrule
\end{tabular}
\end{table}

\subsection{关键区分：自验收 vs 独立验收}

\textbf{Cursor 自验收}：运行 build、启动服务、对照 testing.mdc 自检。Cursor 自己判断自己写的东西能不能用。

\textbf{PM 独立验收}：站在用户角度对照需求清单检查，执行查证三问（目标达成？细节硬伤？换位质疑？）。\textbf{不依赖} Cursor 的自测结论。

\textbf{验收不通过}：PM 记录问题 $\to$ 反馈 Cursor $\to$ Cursor 修复并重新自验收 $\to$ PM 重新独立验收 $\to$ 循环直到通过。

\subsection{交付物}

\begin{itemize}[nosep]
  \item \texttt{AGENT-ARCHITECTURE.md} -- 智能体搭配架构说明（含完整分工矩阵和调用链路）
  \item \texttt{TEAM-MANUAL.md} -- 团队协作说明书（角色分工表、日常场景说明）
\end{itemize}

\subsection{验收结论}

\textcolor{success}{通过。} 分工矩阵已写入 AGENT-ARCHITECTURE.md 和 TEAM-MANUAL.md，自验收与独立验收已明确区分。

% ============================================================
\section{D-003：Skill/规则体系搭建}
% ============================================================

\subsection{需求描述}

为两种智能体建立各自的上下文加载机制。OpenClaw 通过 SKILL.md 获取行为指令，Cursor 通过 MDC 规则文件获取编码约束。两套体系独立运作，但内容需保持一致。

\subsection{核心区别}

\begin{table}[h]
\centering
\begin{tabular}{p{2.5cm}p{4.5cm}p{4.5cm}}
\toprule
\textbf{维度} & \textbf{OpenClaw Skill} & \textbf{Cursor Rule} \\
\midrule
文件格式 & SKILL.md（Markdown） & *.mdc（YAML + Markdown） \\
存放位置 & \texttt{skills/xxx/} & \texttt{.cursor/rules/} \\
加载方式 & AI 按需 read & 自动扫描（alwaysApply / globs） \\
记忆 & 有（memory 文件） & 无（每次全新） \\
\bottomrule
\end{tabular}
\end{table}

\subsection{交付物}

OpenClaw 侧（4 个 Skill）：
\begin{itemize}[nosep]
  \item \texttt{skills/pm/SKILL.md} -- 需求拆解、任务分配、审计、查证
  \item \texttt{skills/tester/SKILL.md} -- 阶段门控验证、功能测试
  \item \texttt{skills/ui-design/SKILL.md} -- 设计方案输出、视觉规范
  \item \texttt{skills/developer/SKILL.md} -- 开发规范（PM 审计用）
\end{itemize}

Cursor 侧（1 + 7 文件）：
\begin{itemize}[nosep]
  \item \texttt{CLAUDE.md} -- 总入口，每次必加载
  \item \texttt{.cursor/rules/project-overview.mdc} -- 技术栈锁定（alwaysApply）
  \item \texttt{.cursor/rules/backend-rules.mdc} -- 后端规范（backend/ 触发）
  \item \texttt{.cursor/rules/frontend-rules.mdc} -- 前端规范（frontend/ 触发）
  \item \texttt{.cursor/rules/ui-design.mdc} -- UI 规范（frontend/ 触发）
  \item \texttt{.cursor/rules/api-conventions.mdc} -- API 规范（backend/ 触发）
  \item \texttt{.cursor/rules/testing.mdc} -- 自测要求（alwaysApply）
  \item \texttt{.cursor/rules/database.mdc} -- 数据库约定（数据库文件触发）
\end{itemize}

\subsection{验收结论}

\textcolor{success}{通过。} 两套体系文件完整，内容无矛盾。

% ============================================================
\section{D-004：文档体系与规范}
% ============================================================

\subsection{需求描述}

建立项目文档的格式规范、命名规则、目录结构和版本管理机制。正式文档统一使用 tex-pdf 格式，确保结构规范性和长期可读性。

\subsection{决策结果}

\textbf{格式}：tex 源文件（Git 管理）+ pdf 交付物（审阅分发），编译工具 Tectonic。

\textbf{命名}：\texttt{\{前缀\}-\{三位编号\}-\{kebab-case 简述\}.tex/pdf}

\begin{table}[h]
\centering
\begin{tabular}{lll}
\toprule
\textbf{前缀} & \textbf{含义} & \textbf{目录} \\
\midrule
REF & 参考文档 & reference/ \\
DEMAND & 需求变更 & demands/ \\
SPEC & 技术规格 & reference/ \\
GUIDE & 操作指南 & reference/ \\
\bottomrule
\end{tabular}
\end{table}

\textbf{目录职责}：
\begin{itemize}[nosep]
  \item \texttt{reference/} -- 参考文档（调研、规范、技术预研），PM AI 主动创建
  \item \texttt{demands/} -- 需求变更文档，每次需求新增/变更/取消生成一份
\end{itemize}

\textbf{版本管理}：语义化版本 vX.Y.Z，每个文档头部标注版本号、日期、变更摘要，CHANGELOG.md 统一记录。

\subsection{交付物}

\begin{itemize}[nosep]
  \item \texttt{reference/DOC-NAMING-CONVENTION.tex/pdf} -- 文档命名与格式规范
  \item \texttt{DOC-LIFECYCLE.md} -- 文档生命周期管理
  \item \texttt{CHANGELOG.md} -- 统一变更日志
\end{itemize}

\subsection{验收结论}

\textcolor{success}{通过。} 规范文档已建立，reference/ 和 demands/ 目录已创建。

% ============================================================
\section{D-005：需求变更流程}
% ============================================================

\subsection{需求描述}

定义需求从提出到交付的完整处理流程，包括记录、评估、设计、拆解、分配、开发、测试、审计、汇报。

\subsection{决策结果}

\textbf{5 步处理流程}：

\begin{enumerate}
  \item \textbf{需求记录}：写入 memory 文件（描述、提出者、优先级）
  \item \textbf{影响评估}（PM 执行）：技术影响 + 文档影响 + 进度影响
  \item \textbf{方案设计}（复杂需求）：更新 SCHEMA.md / SKILL.md，技术栈变更须魏子政审批
  \item \textbf{任务拆解与分配}：拆为可独立完成的开发任务，写入 PROGRESS.md
  \item \textbf{开发--测试--审计--交付}：Cursor 执行 $\to$ PM 审计 $\to$ 文档同步 $\to$ 汇报
\end{enumerate}

\textbf{优先级}：P0 阻塞（立即）、P1 严重（当天）、P2 一般（排入阶段）、P3 轻微（有空再做）。

\textbf{变更控制}：
\begin{itemize}[nosep]
  \item 小变更（不影响已有模块）：PM 直接处理
  \item 中变更（影响已有模块）：PM 评估后执行
  \item 大变更（架构/技术栈/数据库）：PM 评估 $\to$ \textbf{魏子政审批} $\to$ 执行
\end{itemize}

\subsection{交付物}

\begin{itemize}[nosep]
  \item \texttt{DOC-LIFECYCLE.md} -- 文档生命周期管理（含需求处理流程）
  \item \texttt{TEAM-MANUAL.md} -- 日常需求场景说明
\end{itemize}

\subsection{验收结论}

\textcolor{success}{通过。} 5 步流程已写入 DOC-LIFECYCLE.md，优先级和变更控制已定义。

% ============================================================
\section{D-006：自闭环机制}
% ============================================================

\subsection{需求描述}

建立从需求提出到交付完成的全链路闭环机制，确保每个环节不遗漏、不脱节。

\subsection{决策结果}

\textbf{5 个闭环}：

\begin{table}[h]
\centering
\begin{tabular}{clp{5cm}p{3cm}}
\toprule
\textbf{闭环} & \textbf{名称} & \textbf{过程} & \textbf{闭合标志} \\
\midrule
1 & 需求闭环 & 记录$\to$评估$\to$设计$\to$拆解$\to$分配 & PROGRESS.md 有任务 \\
2 & 开发闭环 & 读rules$\to$写代码$\to$自测$\to$提交 & build通过 \\
3 & 测试闭环 & 接口$\to$功能$\to$流程$\to$UI & 验证项全部通过 \\
4 & 文档闭环 & 更新文档$\to$生成PDF$\to$记录CHANGELOG & CHANGELOG有记录 \\
5 & 审计闭环 & 查证三问 & 三问全部通过 \\
\bottomrule
\end{tabular}
\end{table}

\textbf{监控标准}：任何闭环未闭合，任务不能标记为「已完成」。

\subsection{交付物}

\begin{itemize}[nosep]
  \item \texttt{DOC-LIFECYCLE.md} -- 第五章自闭环机制
  \item \texttt{reference/REF-002-project-master-plan.tex/pdf} -- 项目总纲第四部分
\end{itemize}

\subsection{验收结论}

\textcolor{success}{通过。} 5 个闭环定义完整，监控标准明确。

% ============================================================
\section{D-007：开源方案调研}
% ============================================================

\subsection{需求描述}

调研开源 AI Agent 应用商店的实现方案，蒸馏可借鉴的设计思路，为本项目架构提供参考。

\subsection{调研结果}

通过 \texttt{npx skills find} 搜索 skills.sh 平台，安装并分析了 2 个核心方案：

\begin{table}[h]
\centering
\begin{tabular}{p{3cm}p{4cm}p{4.5cm}}
\toprule
\textbf{方案} & \textbf{架构} & \textbf{核心特点} \\
\midrule
Flow Nexus App Store & MCP 工具调用 & 8大分类、模板化部署、AI推荐、质量保障标准 \\
Membrane Marketplacer & 中间件代理 & 免密钥、40+实体模型、连接状态机、MIT开源 \\
\bottomrule
\end{tabular}
\end{table}

\textbf{可借鉴}：分类体系、质量保障标准、模板化部署、免密钥代理思路。\\
\textbf{需自研}：审批工作流、三层安全扫描、企业级 RBAC、审计追溯。\\
\textbf{结论}：现有方案均不完整，本项目架构需自研。

\subsection{交付物}

\begin{itemize}[nosep]
  \item \texttt{reference/REF-001-skill-store-arch-survey.tex/pdf} -- 架构调研报告
\end{itemize}

\subsection{验收结论}

\textcolor{success}{通过。} 两个方案已深入分析，对比结论已记录。

% ============================================================
\section{V0 交付物总表}
% ============================================================

\begin{table}[h]
\centering
\small
\begin{tabular}{p{5.5cm}p{6cm}p{1.8cm}}
\toprule
\textbf{文件} & \textbf{说明} & \textbf{格式} \\
\midrule
\multicolumn{3}{l}{\textbf{tex-pdf 文档（reference/ + demands/）}} \\
\midrule
reference/DOC-NAMING-CONVENTION & 文档命名与格式规范 & tex-pdf \\
reference/REF-001-skill-store-arch-survey & 开源 Skill 商店架构调研 & tex-pdf \\
reference/REF-002-project-master-plan & 项目总纲（V0 核心交付物） & tex-pdf \\
demands/DEMAND-000-v0-master & 本文件：V0 需求总纲 & tex-pdf \\
demands/DEMAND-001-role-redefine & 角色分工重新定义（详细版） & tex-pdf \\
\midrule
\multicolumn{3}{l}{\textbf{Markdown 文档（项目根目录）}} \\
\midrule
AGENT-ARCHITECTURE.md & 智能体搭配架构详细说明 & Markdown \\
tech-selection.md & 技术选型说明书 v2.0 & Markdown \\
TEAM-MANUAL.md & 团队协作说明书 & Markdown \\
DOC-LIFECYCLE.md & 文档生命周期管理 & Markdown \\
CHANGELOG.md & 统一变更日志 & Markdown \\
PROGRESS.md & 进度追踪表 & Markdown \\
BUGS.md & Bug 记录 & Markdown \\
CLAUDE.md & Cursor 总入口 & Markdown \\
\bottomrule
\end{tabular}
\end{table}

% ============================================================
\section{V0 之后的下一步}
% ============================================================

\subsection{V0 待完善项}

\begin{enumerate}[nosep]
  \item 强化 \texttt{ui-design.mdc}，补充常用页面模式（列表页、详情页、表单页、审批页）
  \item 从 cursor.directory 精选 a11y/响应式相关规则片段
  \item 评估 Figma MCP 集成可行性（如魏子政提供设计稿）
  \item 确认数据库 schema 与技术选型的最终一致性
\end{enumerate}

\subsection{V1 阶段里程碑预览}

\begin{table}[h]
\centering
\begin{tabular}{lll}
\toprule
\textbf{阶段} & \textbf{核心交付物} & \textbf{预计周期} \\
\midrule
环境搭建 & Docker Compose + 登录页 + /health & 1 周 \\
核心功能 & 应用 CRUD + 文件上传 + 用户认证 & 1 周 \\
搜索 + 审批 & OpenSearch + 审批流程 & 1 周 \\
安全扫描 & Semgrep + ClamAV + LLM 集成 & 1 周 \\
打磨 + 部署 & UI 优化 + Docker 生产部署 & 2 周 \\
\bottomrule
\end{tabular}
\end{table}

V1 的正式需求将在本阶段完成后，以 DEMAND-002、DEMAND-003 等编号逐条生成，遵循 V0 建立的文档规范和需求处理流程。

% ============================================================
\section*{更新记录}
% ============================================================

\begin{tabular}{llll}
\toprule
版本 & 日期 & 变更内容 & 作者 \\
\midrule
v0.1 & 2026-05-06 & 初始版本：V0 全部 7 项需求整理 & PM AI \\
\bottomrule
\end{tabular}

\end{document}
